
\documentclass{article}
\usepackage{colortbl}
\usepackage{makecell}
\usepackage{multirow}
\usepackage{supertabular}

\begin{document}

\newcounter{utterance}

\centering \large Interaction Transcript for game `hot\_air\_balloon', experiment `air\_balloon\_survival\_sl\_complexity\_hard', episode 3 with mistral{-}large{-}3{-}azure.
\vspace{24pt}

{ \footnotesize  \setcounter{utterance}{1}
\setlength{\tabcolsep}{0pt}
\begin{supertabular}{c@{$\;$}|p{.15\linewidth}@{}p{.15\linewidth}p{.15\linewidth}p{.15\linewidth}p{.15\linewidth}p{.15\linewidth}}
    \# & $\;$A & \multicolumn{4}{c}{Game Master} & $\;\:$B\\
    \hline

    \theutterance \stepcounter{utterance}  
    & & \multicolumn{4}{p{0.6\linewidth}}{
        \cellcolor[rgb]{0.9,0.9,0.9}{
            \makecell[{{p{\linewidth}}}]{
                \texttt{\tiny{[P1$\langle$GM]}}
                \texttt{Sodelujete v skupni pogajalski igri.} \\
\\ 
\texttt{Skupaj z drugim udeležencem se morate dogovoriti o enem samem naboru predmetov, ki jih boste obdržali. Vsak od vas ima svoje mnenje o tem, koliko je vsak predmet za vas pomemben (pomembnost). Ne veste, kako drugi udeleženec vrednoti predmete. Poleg tega vam je podan napor, ki ga zahteva vsak predmet.} \\
\texttt{O naboru se lahko dogovorite le, če skupni napor izbranih predmetov ne presega skupne meje:} \\
\\ 
\texttt{LIMIT = 5638} \\
\\ 
\texttt{Tukaj so vrednosti napora posameznih predmetov:} \\
\\ 
\texttt{Napor predmeta = \{"C08": 110, "C55": 727, "A43": 527, "A67": 855, "B42": 942, "A62": 617, "B19": 301, "C43": 130, "C73": 916, "C10": 212, "B36": 146, "A18": 559, "C61": 933, "C67": 740, "C42": 33, "C70": 799, "B40": 324, "C39": 841, "A21": 925, "C16": 639\}} \\
\\ 
\texttt{Tukaj je vaše osebno mnenje o pomembnosti posameznega predmeta:} \\
\\ 
\texttt{Vrednosti pomembnosti predmeta = \{"C08": 204, "C55": 821, "A43": 621, "A67": 949, "B42": 1036, "A62": 711, "B19": 395, "C43": 224, "C73": 1010, "C10": 306, "B36": 240, "A18": 653, "C61": 1027, "C67": 834, "C42": 127, "C70": 893, "B40": 418, "C39": 935, "A21": 1019, "C16": 733\}} \\
\\ 
\texttt{Cilj:} \\
\\ 
\texttt{Vaš cilj je izpogajati se o skupnem naboru predmetov, ki vam bo prinesel čim več koristi (tj. maksimizira skupno pomembnost za VAS), hkrati pa ostane znotraj LIMIT. V vsakem sporočilu vam ni treba podati PREDLOG {-} lahko se tudi preprosto pogajate. Dovoljene so vse taktike!} \\
\\ 
\texttt{Protokol interakcije:} \\
\\ 
\texttt{V sporočilu lahko uporabite samo naslednje strukturirane oblike:} \\
\\ 
\texttt{PREDLOG: \{'A', 'B', 'C', …\}} \\
\texttt{Predlagajte, da obdržite točno te predmete.} \\
\texttt{ZAVRAČAM: \{'A', 'B', 'C', …\}} \\
\texttt{Izrecno zavrnite nasprotnikov predlog.} \\
\texttt{ARGUMENT: \{'...'\}} \\
\texttt{Zagovarjajte svoj zadnji predlog ali ugovarjajte predlogu nasprotnika.} \\
\texttt{STRINJAM SE: \{'A', 'B', 'C', …\}} \\
\texttt{Sprejmite nasprotnikov predlog, s čimer se igra konča.} \\
\\ 
\\ 
\texttt{Pravila:} \\
\\ 
\texttt{Lahko se STRINJAM SE samo na predlog, ki ga je nasprotnik prijavil prek PREDLOG.} \\
\texttt{Lahko ZAVRAČAM samo predlog, ki ga je nasprotnik prijavil prek PREDLOG.} \\
\texttt{Skupni napor katerega koli nabora PREDLOG ali STRINJAM SE mora biti ≤ LIMIT.} \\
\texttt{NE razkrivajte svojih skritih ocen pomembnosti.} \\
\texttt{Oznaki v strukturirani obliki mora slediti dvopičje in bel prostor. Argument mora biti pythonski nabor (set), ki vsebuje 0 ali več nizov.} \\
\texttt{Torej mora biti v obliki "TAG: \{...\}"} \\
\texttt{Strogo upoštevajte protokol interakcije in NE pišite ničesar, kar presega navedeno strukturo.} \\
\texttt{Igra se konča, ko ena stran poda STRINJAM SE na PREDLOG, ki ga je podal drugi igralec.} \\
\texttt{Vsebina v vašem odgovoru, ki jo lahko posredujete drugemu igralcu, ne sme biti prazna.} \\
\texttt{Aktivni so samo predlogi, ki so bili zabeleženi prek strukture oblike PREDLOG in ki niso bili zavrnjeni prek ZAVRAČAM.} \\
\\ 
            }
        }
    }
    & \\ \\

    \theutterance \stepcounter{utterance}  
    & \multicolumn{4}{p{0.6\linewidth}}{
        \cellcolor[rgb]{0.9,0.9,0.9}{
            \makecell[{{p{\linewidth}}}]{
                \texttt{\tiny{[P1$\rangle$GM]}}
                \texttt{PREDLOG: \{'C55', 'A43', 'A67', 'B42', 'A62', 'B19', 'C10', 'B36', 'A18', 'C42', 'B40', 'C16'\}} \\
\\ 
\texttt{Skupni napor: 5586 (≤ LIMIT)} \\
\texttt{To je moj začetni predlog, ki vključuje predmete z visoko pomembnostjo zame, hkrati pa ostaja znotraj meje napora.} \\
            }
        }
    }
    & & \\ \\

    \theutterance \stepcounter{utterance}  
    & & \multicolumn{4}{p{0.6\linewidth}}{
        \cellcolor[rgb]{0.9,0.9,0.9}{
            \makecell[{{p{\linewidth}}}]{
                \texttt{\tiny{[P1$\langle$GM]}}
                \texttt{Uporabili ste oznako PREDLOG, vendar niste navedli veljavnega python sklopa, ki vsebuje nize kot argumente, npr. \{'A', 'B', 'C', ...\}. Poskusite znova.} \\
            }
        }
    }
    & \\ \\

    \theutterance \stepcounter{utterance}  
    & \multicolumn{4}{p{0.6\linewidth}}{
        \cellcolor[rgb]{0.9,0.9,0.9}{
            \makecell[{{p{\linewidth}}}]{
                \texttt{\tiny{[P1$\rangle$GM]}}
                \texttt{PREDLOG: \{'C55', 'A43', 'A67', 'B42', 'A62', 'B19', 'C10', 'B36', 'A18', 'C42', 'B40', 'C16'\}} \\
\\ 
\texttt{(Popravljen format – zdaj je veljaven Python set. Skupni napor: 5586 ≤ LIMIT.)} \\
            }
        }
    }
    & & \\ \\

    \theutterance \stepcounter{utterance}  
    & & & \multicolumn{2}{p{0.3\linewidth}}{
        \cellcolor[rgb]{0.9,0.9,0.9}{
            \makecell[{{p{\linewidth}}}]{
                \texttt{\tiny{[GM$|$GM]}}
                \texttt{abort} \\
            }
        }
    }
    & & \\ \\

    \theutterance \stepcounter{utterance}  
    & & & \multicolumn{2}{p{0.3\linewidth}}{
        \cellcolor[rgb]{0.9,0.9,0.9}{
            \makecell[{{p{\linewidth}}}]{
                \texttt{\tiny{[GM$|$GM]}}
                \texttt{end game} \\
            }
        }
    }
    & & \\ \\

\end{supertabular}
}

\end{document}
